This chapter provides a self-contained summary of the thesis. Section 1.1 introduces the context and topic. Section 1.2 summarizes the state of the art and identifies the research gap. Section 1.3 states the research questions. Section 1.4 outlines the method and approach. Section 1.5 presents the findings. The subsequent chapters provide detailed treatment: Chapter 2 establishes the background, Chapter 3 formulates the research questions and methodology, Chapter 4 presents the literature review, Chapter 5 develops the conceptual framework, Chapter 6 describes the prototype implementation, Chapter 7 reports the evaluation, and Chapter 8 concludes the thesis.

\section{Context and Topic}

Industrial automation, energy systems, transportation infrastructures, and manufacturing plants rely on continuous and dependable network connectivity. Industrial routers function as the communication backbone in these environments, enabling remote access, secure data exchange, device management, and operational monitoring. Interruptions in network equipment functionality whether caused by software faults, hardware defects, power supply failures, or environmental interference lead to halted production processes, violations of safety requirements, and operational losses. The need for predictable behavior, determinism, and reliability establishes fault management as a central requirement for industrial networking equipment.

Modern industrial router platforms increasingly follow \gls{soa} principles, where routing, \gls{vpn} management, firewalling, device configuration, and supervisory control are executed as modular and interacting services. In this context, a platform refers to the combined hardware and software environment on which these services execute, including the operating system, message bus, configuration engine, and service lifecycle manager. While service modularity promotes maintainability and independent evolution of components, it also introduces failure modes that do not exist in monolithic software: faults originating in one service propagate along service dependency paths through inter-service communication channels and shared state, resulting in cascading failures or degraded performance in dependent services. This propagation behavior is a direct consequence of the service decomposition monolithic systems confine faults within process boundaries, whereas SOA systems expose them across service interfaces. Critical functions such as routing and VPN termination may employ redundancy mechanisms, but auxiliary services (logging, certificate management, telemetry export) typically do not, making them vulnerable to single points of failure. Existing fault taxonomies structured classification schemes that categorize faults by type, origin, severity, and recoverability provide useful classifications for service-oriented and distributed systems, but they do not define operational models for managing faults in resource-constrained environments such as industrial routers.

\section{State of the Art}

Research in distributed systems offers advanced techniques for observability, anomaly detection, and root cause localization, often relying on large-scale telemetry sets, machine learning methods, and high-performance computing resources. These approaches are designed for cloud-native microservice environments and are generally not directly applicable to embedded industrial router platforms that require deterministic operation, transparent fault reasoning, and bounded computational overhead. The term "deterministic" is used throughout this thesis to refer to four specific properties defined in the Definitions section: deterministic decision logic, deterministic processing order, deterministic recovery sequencing, and deterministic network behavior ensuring that fault management activities do not disrupt the forwarding plane the router's data-path responsible for packet forwarding, routing, and network address translation. Self-healing frameworks for SOA-based systems additionally focus on orchestration-heavy environments and often do not incorporate explicit, formalized episode semantics that enable bounded, deterministic recovery behavior suitable for long-lived industrial routers.

The current state of research therefore shows a lack of a unified, stateful framework that integrates structured telemetry organization, explicit modelling of service dependencies and failure propagation, and deterministic recovery strategies suitable for industrial-grade SOA-based routers. In particular, existing approaches rarely provide explicit runtime episode semantics that distinguish new faults from ongoing violations, ensuring that recovery actions execute at most once per fault episode. They also seldom combine root-cause diagnosis with impact-aware suppression of dependent symptoms in order to prevent alert storms and redundant recovery actions. This gap results in fragmented fault management practices, limited visibility into service interdependencies, and inconsistent recovery behavior across router platforms. A detailed analysis of the literature is provided in Chapter 4.

\section{Research Questions}

The central problem addressed in this thesis is the absence of a stateful model for fault detection, isolation, and recovery (\gls{fdir}) in SOA-based industrial routers. Five research questions guide the investigation: (RQ1) how faults can be systematically characterized for detection, isolation, and recovery under embedded constraints; (RQ2) how telemetry can be organized to enable stateful fault tracking and root cause localization; (RQ3) how service dependencies should be represented to enable deterministic reasoning about failure propagation; (RQ4) which deterministic recovery strategies are suitable and how episode-aware coordination can prevent redundant actions; and (RQ5) how these elements can be integrated into a unified stateful FDIR architecture. The detailed formulation of these questions is provided in Chapter 3.

\section{Method and Approach}

The purpose of the thesis is to define a stateful FDIR model by synthesizing existing research in fault detection, telemetry, dependency analysis, and deterministic recovery, and by adapting these concepts to the operational constraints of industrial router environments. The methodology follows a structured research process comprising eight steps: a systematic literature review across five domains (Section 3.3), identification of relevant fault types (RQ1), design of a structured telemetry model (RQ2), modelling of service dependencies (RQ3), identification of deterministic recovery strategies (RQ4), integration into a stateful FDIR framework (RQ5), scenario-based evaluation, and validation of research questions against empirical evidence. The conceptual framework is presented in Chapter 5. A prototype implementation (prism-police) on the Belden/Hirschmann PRISM industrial router platform realizes the framework as a C++20 application (Chapter 6). The evaluation (Chapter 7) validates the prototype through 47 unit tests and quantitative experiments comparing the prototype with Prometheus and Nagios.

\section{Findings}

The framework is designed to be platform-independent in its conceptual architecture all behavior is defined through declarative configuration rather than hardcoded service logic though the prototype implementation and evaluation are conducted on a single platform (Belden/Hirschmann PRISM). Quantitative single-run experiments demonstrate that the episode model reduces alert volume by two to five orders of magnitude under sustained fault conditions, while CPU utilization remains at 77.49\% of a single core versus 313.83\% (Prometheus) and 1,035.97\% (Nagios) at 1,000 samples per second. Memory consumption stays between 13.92 and 29.37 MB across all tested scenarios. The at-most-once recovery execution guarantee is verified through unit tests. The evaluation is limited to a single platform and desktop-class hardware; production-platform validation, long-duration resource profiling, and cross-platform applicability remain open. Validation of cross-platform applicability remains future work. The following chapters build on this introduction by first establishing the relevant technological foundations and then analyzing the existing literature through the lens of runtime fault episode management, deterministic recovery coordination, and embedded system constraints.